\textbf{Phần 1. Trắc nghiệm lựa chọn}
\setcounterpageref{ex}{0}

\begin{ex} %[3D11N61]
Làm tròn số $4\,320$ đến hàng nghìn ta được
\choice
{$5000$}
{$4300$}
{$4500$}
{\True $4000$}
\loigiai{
Chữ số hàng trăm là 3 ($3 < 5$) nên làm tròn xuống $4000$.\\
\textbf{Đáp án của câu hỏi là D}
}
\end{ex}

\begin{ex} %[3D11N61]
Làm tròn số $7\,850$ đến hàng nghìn ta được
\choice
{$7000$}
{$7800$}
{\True $8000$}
{$9000$}
\loigiai{
Hàng trăm là 8 ($\ge 5$) nên làm tròn lên $8000$.\\
\textbf{Đáp án của câu hỏi là C}
}
\end{ex}

\begin{ex} %[3D11H61]
Làm tròn số $34\,500$ đến hàng nghìn ta được
\choice
{$34000$}
{$30000$}
{$40000$}
{\True $35000$}
\loigiai{
Hàng trăm là 5 ($\ge 5$) nên làm tròn lên $35000$.\\
\textbf{Đáp án của câu hỏi là D}
}
\end{ex}

\begin{ex} %[3D11H61]
Làm tròn số $62\,300$ đến hàng chục nghìn ta được
\choice
{$70000$}
{\True $60000$}
{$62000$}
{$65000$}
\loigiai{
Hàng nghìn là 2 ($< 5$) nên làm tròn thành $60000$.\\
\textbf{Đáp án của câu hỏi là B}
}
\end{ex}

\begin{ex} %[3D11H61]
Làm tròn số $87\,400$ đến hàng chục nghìn ta được
\choice
{\True $90000$}
{$80000$}
{$87000$}
{$100000$}
\loigiai{
Hàng nghìn là 7 ($\ge 5$) nên làm tròn lên $90000$.\\
\textbf{Đáp án của câu hỏi là A}
}
\end{ex}

\begin{ex} %[3D11H61]
Làm tròn số $45\,000$ đến hàng chục nghìn ta được
\choice
{$40000$}
{\True $50000$}
{$45000$}
{$60000$}
\loigiai{
Hàng nghìn là 5 ($\ge 5$) nên làm tròn lên $50000$.\\
\textbf{Đáp án của câu hỏi là B}
}
\end{ex}

\begin{ex} %[3D11V61]
Một quãng đường dài $12\,780\text{ m}$. Làm tròn đến hàng nghìn ta được khoảng
\choice
{$12000\text{ m}$}
{$12800\text{ m}$}
{\True $13000\text{ m}$}
{$10000\text{ m}$}
\loigiai{
Hàng trăm là 7 ($\ge 5$) nên làm tròn lên $13000\text{ m}$.\\
\textbf{Đáp án của câu hỏi là C}
}
\end{ex}

\begin{ex} %[3D11V61]
Một sân vận động có $43\,250$ khán giả. Làm tròn đến hàng chục nghìn ta được khoảng
\choice
{$50000$ khán giả}
{$43000$ khán giả}
{$45000$ khán giả}
{\True $40000$ khán giả}
\loigiai{
Hàng nghìn là 3 ($< 5$) nên làm tròn thành $40000$ khán giả.\\
\textbf{Đáp án của câu hỏi là D}
}
\end{ex}

\begin{ex} %[3D11V61]
Dân số một thị trấn là $28\,650$ người. Làm tròn đến hàng nghìn ta được khoảng
\choice
{$28000$ người}
{$30000$ người}
{\True $29000$ người}
{$28600$ người}
\loigiai{
Hàng trăm là 6 ($\ge 5$) nên làm tròn lên $29000$ người.\\
\textbf{Đáp án của câu hỏi là C}
}
\end{ex}

\begin{ex} %[3D11V61]
Làm tròn số $75\,800$ đến hàng chục nghìn ta được
\choice
{\True $80000$}
{$70000$}
{$76000$}
{$90000$}
\loigiai{
Hàng nghìn là 5 nên làm tròn lên $80000$.\\
\textbf{Đáp án của câu hỏi là A}
}
\end{ex}

\begin{ex} %[3D11V61]
Có bao nhiêu số tròn trăm khi làm tròn đến hàng nghìn được số $6\,000$?
\choice
{$9$ số}
{\True $10$ số (từ $5500$ đến $6400$)}
{$11$ số}
{$8$ số}
\loigiai{
Các số tròn trăm từ 5500 đến 6400 gồm có 10 số.\\
\textbf{Đáp án của câu hỏi là B}
}
\end{ex}

\begin{ex} %[3D11C61]
Tìm số lớn nhất có 5 chữ số mà khi làm tròn đến hàng chục nghìn thì được số $70\,000$.
\choice
{\True $74999$}
{$75000$}
{$74000$}
{$69999$}
\loigiai{
Để làm tròn thành 70000 thì chữ số hàng nghìn phải $< 5$. Số lớn nhất là $74999$.\\
\textbf{Đáp án của câu hỏi là A}
}
\end{ex}

\textbf{Phần 2. Trắc nghiệm đúng sai}
\setcounterpageref{ex}{0}

\begin{ex} %[3D11H61]
Xét tính đúng sai của việc làm tròn số:
\choiceTF
{\True Làm tròn $34\,200$ đến hàng nghìn được $34\,000$.}
{\True Làm tròn $78\,500$ đến hàng chục nghìn được $80\,000$.}
{Làm tròn $54\,999$ đến hàng chục nghìn được $60\,000$.}
{\True Làm tròn $54\,999$ đến hàng chục nghìn được $50\,000$ (vì hàng nghìn là $4 < 5$).}
\loigiai{
\begin{itemize}
    \item $34200 \rightarrow 34000$: Đúng.
    \item $78500 \rightarrow 80000$: Đúng.
    \item Hàng nghìn là 4 nên làm tròn thành 50000. Ý 3 Sai, Ý 4 Đúng.
\end{itemize}
\begin{itemize}
    \item a) \textbf{Đúng}.
    \item b) \textbf{Đúng}.
    \item c) \textbf{Sai}.
    \item d) \textbf{Đúng}.
\end{itemize}
}
\end{ex}

\begin{ex} %[3D11V61]
Cho số $65\,480$. Xét tính đúng sai:
\choiceTF
{\True Làm tròn đến hàng nghìn được $65\,000$.}
{\True Làm tròn đến hàng chục nghìn được $70\,000$.}
{Làm tròn đến hàng nghìn được $66\,000$.}
{\True Chữ số hàng trăm là $4 < 5$ nên giữ nguyên hàng nghìn.}
\loigiai{
\begin{itemize}
    \item Hàng nghìn: $65000$. Đúng.
    \item Hàng chục nghìn: $70000$ (vì hàng nghìn là $5 \ge 5$). Đúng.
    \item Hàng nghìn được 65000 (không phải 66000). Sai.
    \item Giải thích đúng: Đúng.
\end{itemize}
\begin{itemize}
    \item a) \textbf{Đúng}.
    \item b) \textbf{Đúng}.
    \item c) \textbf{Sai}.
    \item d) \textbf{Đúng}.
\end{itemize}
}
\end{ex}

\textbf{Phần 3. Trả lời ngắn}
\setcounterpageref{ex}{0}

\begin{ex} %[3D11H61]
Làm tròn số $8\,650$ đến hàng nghìn được số nào?
\shortans{9000}
\loigiai{
\textbf{Đáp số:} $9000$\\
$9000$.
}
\end{ex}

\begin{ex} %[3D11H61]
Làm tròn số $42\,000$ đến hàng chục nghìn được số nào?
\shortans{40000}
\loigiai{
\textbf{Đáp số:} $40000$\\
$40000$.
}
\end{ex}

\begin{ex} %[3D11V61]
Làm tròn số $57\,340$ đến hàng nghìn được số nào?
\shortans{57000}
\loigiai{
\textbf{Đáp số:} $57000$\\
$57000$.
}
\end{ex}

\begin{ex} %[3D11V61]
Làm tròn số $85\,400$ đến hàng chục nghìn được số nào?
\shortans{90000}
\loigiai{
\textbf{Đáp số:} $90000$\\
$90000$.
}
\end{ex}

\begin{ex} %[3D11V61]
Tìm số bé nhất có 5 chữ số mà khi làm tròn đến hàng chục nghìn thì được số $60\,000$.
\shortans{55000}
\loigiai{
\textbf{Đáp số:} $55000$\\
$55000$.
}
\end{ex}

\begin{ex} %[3D11C61]
Có bao nhiêu số tròn nghìn khi làm tròn đến hàng chục nghìn thì được số $40\,000$?
\shortans{10}
\loigiai{
\textbf{Đáp số:} $10$\\
Các số tròn nghìn từ 35000 đến 44000 gồm có 10 số.
}
\end{ex}
